\documentclass{amsart}
\usepackage{amsmath,amssymb,amsthm,hyperref}
\usepackage{algorithm}
\usepackage{algpseudocode}

\newtheorem{theorem}{Theorem}
\newtheorem{lemma}{Lemma}
\newtheorem{proposition}{Proposition}
\newtheorem{corollary}{Corollary}
\theoremstyle{definition}
\newtheorem{definition}{Definition}
\newtheorem{remark}{Remark}

\newcommand{\asymp}{\sim}

\begin{document}

\title{Asymptotics of the interface width for an expanding Family--Vicsek interface}
\author{}
\date{}
\maketitle

\section{Formalization and standing assumptions}

We make precise the (semi-empirical) scaling hypotheses of the problem, state them
as axioms, and derive everything rigorously from those axioms. The mathematical
content is the analysis of a one-parameter family of competing power laws; the
experimental content is the protocol of Section~\ref{sec:protocol}.

Throughout, ``$g(t)\asymp h(t)$ as $t\to\infty$'' means there are constants
$0<c\le C<\infty$ and $t_0$ with $c\,h(t)\le g(t)\le C\,h(t)$ for all $t\ge t_0$;
for two power laws this is equivalent to equality of exponents (with the prefactor
absorbed into $c,C$), which is what we track. We write $f(u)\asymp k(u)$ as
$u\to0^+$ or $u\to\infty$ analogously.

\begin{definition}[Family--Vicsek scaling on a fixed substrate]\label{def:FV}
Fix exponents $\alpha>0$ (roughness) and $\beta>0$ (growth), and set
$z:=\alpha/\beta>0$ (dynamic exponent). On a substrate of fixed linear size
$L_0$ the height field $h(x,t)$ has interface variance
\begin{equation}\label{eq:FV}
  W_0^2(L_0,t):=\langle h(x,t)^2\rangle-\langle h(x,t)\rangle^2
  \;\asymp\; L_0^{2\alpha}\, f\!\left(\frac{t}{L_0^{\,z}}\right),
\end{equation}
where the dimensionless positive scaling function $f$ obeys
\begin{equation}\label{eq:fasymp}
  f(u)\asymp u^{2\beta}\ \ (u\to 0^+),\qquad
  f(u)\to f_\infty\in(0,\infty)\ \ (u\to\infty),
\end{equation}
and is continuous and strictly positive on $(0,\infty)$, so that
$0<\inf_{[a,b]}f\le\sup_{[a,b]}f<\infty$ on every compact $[a,b]\subset(0,\infty)$.
\end{definition}

\begin{lemma}[Growth-phase size independence]\label{lem:sizeindep}
On a fixed substrate, before saturation (i.e.\ for $t\ll L_0^{z}$) the
Family--Vicsek width is independent of the substrate size:
$W_0^2(L_0,t)\asymp t^{2\beta}$, with implicit constant independent of $L_0$.
After saturation ($t\gg L_0^{z}$) it is time-independent and size-dependent:
$W_0^2(L_0,t)\asymp f_\infty\,L_0^{2\alpha}$.
\end{lemma}
\begin{proof}
For $t\ll L_0^{z}$, $u:=t/L_0^{z}\to0$, so by \eqref{eq:fasymp}
$f(u)\asymp u^{2\beta}$, hence by \eqref{eq:FV}
$W_0^2\asymp L_0^{2\alpha}(t/L_0^{z})^{2\beta}=t^{2\beta}\,L_0^{\,2\alpha-2\beta z}
=t^{2\beta}$, because $2\alpha-2\beta z=2\alpha-2\beta(\alpha/\beta)=0$. For
$t\gg L_0^{z}$, $u\to\infty$, so $f(u)\to f_\infty$ and
$W_0^2\asymp f_\infty L_0^{2\alpha}$.
\end{proof}

Lemma~\ref{lem:sizeindep} is the structural hallmark of Family--Vicsek scaling and
will be decisive for the extraction analysis: \emph{the substrate size enters the
width only after saturation}.

The single dimensionless control variable in \eqref{eq:FV} is the ratio of lateral
correlation length to system size. Defining the \emph{correlation length}
\begin{equation}\label{eq:xi}
  \xi(t):=t^{1/z},
\end{equation}
the argument of $f$ is $t/L_0^{z}=(\xi(t)/L_0)^{z}$, with crossover at
$\xi(t)=L_0$.

\begin{definition}[Expanding substrate; adiabatic Family--Vicsek closure]\label{def:expand}
The substrate grows as $L(t)\asymp t^{\gamma}$ ($\gamma>0$) for $t\ge t_0$, with
reference size $L_0:=L(t_0)$ at a fixed reference time $t_0>0$; thus
\begin{equation}\label{eq:Lt}
  L(t)=L_0\,(t/t_0)^{\gamma},\qquad t\ge t_0,\ \gamma>0.
\end{equation}
We posit the \emph{adiabatic Family--Vicsek closure}: the expanding-interface
variance obeys \eqref{eq:FV}--\eqref{eq:fasymp} with the fixed size replaced at
each instant by the current size $L(t)$:
\begin{equation}\label{eq:adiabatic}
  W^2(t):=\langle h(x,t)^2\rangle-\langle h(x,t)\rangle^2
  \;\asymp\; L(t)^{2\alpha}\, f\!\left(\frac{t}{L(t)^{\,z}}\right)
  = L(t)^{2\alpha}\, f\!\left(\Big(\tfrac{\xi(t)}{L(t)}\Big)^{z}\right),
\end{equation}
with the same $\alpha,\beta,z,f$ as in Definition~\ref{def:FV}.
\end{definition}

\begin{remark}[Status of the closure]
Equation~\eqref{eq:adiabatic} is the standard closure for slowly expanding
interfaces: on time scales over which $L(t)$ changes by an $O(1)$ factor, the
interface relaxes onto the Family--Vicsek form for the instantaneous size. It is
the hypothesis the problem invites us to analyze; all results below are theorems
\emph{conditional on \eqref{eq:adiabatic}}.
\end{remark}

\section{The dichotomy theorem: exactly two leading regimes (plus the boundary)}
\label{sec:dichotomy}

\begin{theorem}[Trichotomy of asymptotics]\label{thm:main}
Assume \eqref{eq:FV}--\eqref{eq:adiabatic}. As $t\to\infty$ the expanding-interface
variance satisfies, with $L(t)=L_0(t/t_0)^{\gamma}$,
\begin{align}
  \text{\emph{(F)} fast/uncorrelated, } \gamma>1/z:\quad
    & W^2(t)\;\asymp\; t^{2\beta}
      \;=\; L(t)^{2\alpha}\,t^{\,2\beta-2\gamma\alpha},
      \label{eq:fast}\\[2pt]
  \text{\emph{(S)} slow/correlated, } \gamma<1/z:\quad
    & W^2(t)\;\asymp\; f_\infty\,L(t)^{2\alpha}
      \;\asymp\; L_0^{2\alpha}\,t^{2\gamma\alpha},
      \label{eq:slow}\\[2pt]
  \text{\emph{(C)} critical, } \gamma=1/z:\quad
    & W^2(t)\;\asymp\; t^{2\beta}\;\asymp\;L(t)^{2\alpha}.
      \label{eq:crit}
\end{align}
In particular $W^2(t)\asymp \kappa(L_0)\,t^{\theta(\gamma)}$ for a single power
\begin{equation}\label{eq:theta}
  \theta(\gamma)=\min\{\,2\beta,\;2\gamma\alpha\,\}
  =\begin{cases}2\beta, & \gamma\ge 1/z,\\[2pt] 2\gamma\alpha, & \gamma\le 1/z,\end{cases}
\end{equation}
and there are \emph{no other} asymptotic power laws. The size-dependence of the
amplitude is
\begin{equation}\label{eq:kappa}
  \kappa(L_0)\asymp
  \begin{cases}
    L_0^{0}\ \ (\text{$L_0$-independent}), & \gamma>1/z,\\[2pt]
    L_0^{2\alpha}\,t_0^{-2\gamma\alpha}\ \ (\text{$\asymp L_0^{2\alpha}$}), & \gamma<1/z,\\[2pt]
    L_0^{0}\ \ (\text{$L_0$-independent}), & \gamma=1/z.
  \end{cases}
\end{equation}
The exponent $\theta(\gamma)$ is continuous, non-decreasing, concave (piecewise
linear with a single kink at $\gamma=1/z$), and equals $2\beta$ on $[1/z,\infty)$.
\end{theorem}

\begin{proof}
Put $\rho(t):=t/L(t)^{z}=(\xi(t)/L(t))^{z}$. By \eqref{eq:Lt},
\begin{equation}\label{eq:rho}
  \rho(t)=\frac{t}{(L_0 (t/t_0)^{\gamma})^{z}}
  =L_0^{-z}t_0^{\gamma z}\,t^{\,1-\gamma z}
  \asymp t^{\,z(1/z-\gamma)} ,
\end{equation}
so $\operatorname{sgn}$ of the time-exponent of $\rho$ equals
$\operatorname{sgn}(1/z-\gamma)$ (as $z>0$).

\medskip\noindent\textbf{Case (F): $\gamma>1/z$.} Then $\rho(t)\to0^+$. By
\eqref{eq:fasymp} there are $0<c_1\le C_1<\infty$, $u_0>0$ with
$c_1u^{2\beta}\le f(u)\le C_1u^{2\beta}$ for $0<u\le u_0$, and some $T$ with
$\rho(t)\le u_0$ for $t\ge T$; hence $f(\rho(t))\asymp\rho(t)^{2\beta}$.
Substituting into \eqref{eq:adiabatic},
\[
  W^2(t)\asymp L(t)^{2\alpha}\rho(t)^{2\beta}
  =L(t)^{2\alpha}\Big(\tfrac{t}{L(t)^{z}}\Big)^{2\beta}
  =t^{2\beta}\,L(t)^{2\alpha-2\beta z}=t^{2\beta},
\]
since $2\alpha-2\beta z=0$. This is \eqref{eq:fast}; the rewrite
$t^{2\beta}=L(t)^{2\alpha}t^{2\beta-2\gamma\alpha}$ uses
$L(t)^{2\alpha}=L_0^{2\alpha}t_0^{-2\gamma\alpha}t^{2\gamma\alpha}$. By
Lemma~\ref{lem:sizeindep} the amplitude is $L_0$-independent, giving the first line
of \eqref{eq:kappa}. As $\gamma>1/z=\beta/\alpha$ gives $2\gamma\alpha>2\beta$, the
exponent is $2\beta=\min\{2\beta,2\gamma\alpha\}$, matching \eqref{eq:theta}.

\medskip\noindent\textbf{Case (S): $\gamma<1/z$.} Then $\rho(t)\to+\infty$, so by
\eqref{eq:fasymp} $f(\rho(t))\to f_\infty\in(0,\infty)$, i.e.\ $f(\rho(t))\asymp1$.
Hence by \eqref{eq:adiabatic} and \eqref{eq:Lt},
\[
  W^2(t)\asymp f_\infty L(t)^{2\alpha}
  =f_\infty L_0^{2\alpha}t_0^{-2\gamma\alpha}\,t^{2\gamma\alpha}
  \asymp L_0^{2\alpha}\,t^{2\gamma\alpha},
\]
which is \eqref{eq:slow}; the amplitude is $\asymp L_0^{2\alpha}$, the second line
of \eqref{eq:kappa}. As $\gamma<\beta/\alpha$ gives $2\gamma\alpha<2\beta$, the
exponent is $2\gamma\alpha=\min\{2\beta,2\gamma\alpha\}$.

\medskip\noindent\textbf{Case (C): $\gamma=1/z$.} Then by \eqref{eq:rho}
$\rho(t)\asymp t^{0}\asymp1$, so $\rho(t)\in[a,b]$ for some $0<a\le b<\infty$ and
all large $t$. By continuity and positivity of $f$ on $[a,b]$
(Definition~\ref{def:FV}), $f(\rho(t))\asymp1$, and
\[
  W^2(t)\asymp L(t)^{2\alpha}\asymp t^{2\gamma\alpha}=t^{2\alpha/z}=t^{2\beta},
\]
using $\gamma=1/z$, $z=\alpha/\beta$; the amplitude $L(t)^{2\alpha}\asymp t^{2\beta}$
carries the explicit $t_0,L_0$ inside the constant, so as a function of $L_0$ at
fixed $t$ the leading $t$-power is $L_0$-independent at the level of the time
exponent. This is \eqref{eq:crit} and the third line of \eqref{eq:kappa}.

\medskip\noindent\textbf{No other behaviour; shape of $\theta$.} The cases
partition $(0,\infty)$, in each $W^2$ is a single power $t^{\theta(\gamma)}$ with
\eqref{eq:theta}. Continuity at $\gamma=1/z$: both branches give
$2\beta=2(1/z)\alpha$. On $(0,1/z]$, $\theta=2\alpha\gamma$ strictly increases
(slope $2\alpha$); on $[1/z,\infty)$, $\theta\equiv2\beta$. Being the minimum of
two affine functions of $\gamma$, $\theta$ is concave with a single downward kink
at $\gamma=1/z$ (slope drops $2\alpha\to0$). This disproves any third regime: the
only freedom in $\theta$ is the kink location $1/z$ and the slopes $2\alpha,0$.
\end{proof}

\begin{remark}[Correction of the problem's first fast-form]\label{rem:correction}
The problem states the fast regime as
``$W^2\asymp L_0^{2\alpha}t^{2\beta}\asymp L(t)^{2\alpha}t^{2\beta-2\gamma\alpha}$.''
The \emph{second} expression is correct (and equals $t^{2\beta}$, our
\eqref{eq:fast}), but the \emph{first} expression $L_0^{2\alpha}t^{2\beta}$ is
\emph{not} consistent with Family--Vicsek scaling: by
Lemma~\ref{lem:sizeindep} the growing (unsaturated) width is \emph{always}
substrate-size independent, because the would-be amplitude
$L_0^{2\alpha}(t/L_0^{z})^{2\beta}$ has $L_0$-exponent $2\alpha-2\beta z=0$. Hence
in the fast regime $W^2\asymp t^{2\beta}$ with an $L_0$-independent constant; the
factor $L_0^{2\alpha}$ in the problem's first form is spurious (it is cancelled by
$L(t)^{-2\alpha}$ hidden in $t^{2\beta-2\gamma\alpha}$ when one rewrites via
$L(t)$). This correction is essential for the extraction analysis: $\alpha$
\emph{cannot} be read off from the $L_0$-dependence of the fast regime, contrary to
the literal phrasing of part (a). It is the slow (saturated) regime that carries
$\alpha$, cf.\ \eqref{eq:kappa}. We prove the corrected extraction in
Section~\ref{sec:protocol}.
\end{remark}

\begin{remark}[The boundary case]
The two strict regimes leave $\gamma=1/z$ unaddressed; there $\rho(t)\asymp1$ so
$f$ sits in its crossover, a genuinely distinct sub-leading structure. But the
leading power $t^{2\beta}$ at $\gamma=1/z$ matches both one-sided limits, so
$\theta(\gamma)$ is continuous and the two strict regimes are the only two
\emph{leading} power laws, separated by the single critical point $\gamma=1/z$.
\end{remark}

\section{The order parameter $1/z-\gamma$: distance from correlation}
\label{sec:correlation}

\begin{definition}[Global correlation]\label{def:corr}
At time $t$ the interface is \emph{globally correlated} (saturated) if
$\xi(t)\ge L(t)$ and \emph{uncorrelated} (roughening) if $\xi(t)<L(t)$.
\end{definition}

\begin{proposition}[Sign law]\label{prop:sign}
Let $\Delta:=1/z-\gamma$. For large $t$,
\[
  \frac{\xi(t)}{L(t)}\asymp t^{\Delta},\qquad
  \frac{\log(\xi(t)/L(t))}{\log t}\to\Delta .
\]
Hence: (i) $\Delta>0\iff\gamma<1/z\iff\xi/L\to\infty$: correlated phase (S);
(ii) $\Delta<0\iff\gamma>1/z\iff\xi/L\to0$: uncorrelated phase (F);
(iii) $\Delta=0\iff\gamma=1/z\iff\xi/L\asymp1$: critical phase (C).
The magnitude $|\Delta|$ is the per-decade rate of approach to/retreat from full
correlation: $\xi/L$ changes by the factor $10^{\Delta}$ when $t\mapsto10t$. Thus
$\Delta=1/z-\gamma$ is the precise ``distance from correlation,'' its sign the
phase and $|\Delta|$ the speed.
\end{proposition}
\begin{proof}
By \eqref{eq:xi},\eqref{eq:Lt}, $\xi/L=t^{1/z}/(L_0(t/t_0)^{\gamma})
\asymp t^{1/z-\gamma}=t^{\Delta}$, giving the limit. The equivalences are the sign
cases of $\Delta$, matched to the cases of Theorem~\ref{thm:main} via
\eqref{eq:rho} ($\rho=(\xi/L)^z\asymp t^{z\Delta}$, same sign). Per decade:
$(\xi/L)(10t)\big/(\xi/L)(t)\asymp10^{\Delta}$.
\end{proof}

\begin{theorem}[Minimal control strength to switch phase]\label{thm:control}
Let the bare system have growth index $\gamma_0$ and dynamic exponent $z$, hence
phase $\operatorname{sgn}(\Delta_0)$ with $\Delta_0:=1/z-\gamma_0$. A control
protocol replacing $\gamma_0$ by $\gamma=\gamma_0+\delta$ ($\delta\in\mathbb R$,
$\gamma>0$) crosses the phase boundary iff $\gamma$ crosses $1/z$, i.e.\ iff
$\delta$ crosses $\Delta_0$. The boundary is reached exactly at $\delta=\Delta_0$
(critical phase (C)), and the system lands strictly in the opposite phase for any
$\delta$ on the far side of $\Delta_0$; the infimal switching strength is
\begin{equation}\label{eq:thresh}
  |\delta_{\min}|=\big|\,\tfrac1z-\gamma_0\,\big|=|\Delta_0|,
\end{equation}
attained in the limit. Directionally: from the uncorrelated phase
($\gamma_0>1/z$) one must \emph{decrease} $\gamma$ with
$\delta<\Delta_0<0$, $|\delta|>|\Delta_0|$; from the correlated phase
($\gamma_0<1/z$) one must \emph{increase} $\gamma$ with $\delta>\Delta_0>0$.
\end{theorem}
\begin{proof}
By Theorem~\ref{thm:main} and Proposition~\ref{prop:sign}, the phase is
$\operatorname{sgn}(1/z-\gamma)=\operatorname{sgn}(\Delta)$, boundary $\Delta=0$.
With $\gamma=\gamma_0+\delta$, $\Delta(\delta)=\Delta_0-\delta$ is affine,
strictly decreasing, with unique zero $\delta=\Delta_0$. Thus the sign flips
exactly as $\delta$ passes $\Delta_0$: same side $\Rightarrow$ same phase, far side
$\Rightarrow$ opposite phase, $\delta=\Delta_0\Rightarrow$ critical. The infimal
$|\delta|$ that changes the sign is $|\Delta_0|$ (not attained at the strict level,
attained as the boundary). The directional statements are the two sign cases of
$\Delta_0$.
\end{proof}

\section{Experimental protocol and unambiguous extraction of $\alpha,\beta$}
\label{sec:protocol}

By Remark~\ref{rem:correction} and \eqref{eq:kappa}, the two exponents live in
\emph{different} regimes:
\begin{itemize}
\item the \textbf{fast} regime (F) yields $\beta$ via the time exponent
  $W^2\asymp t^{2\beta}$ (and is $L_0$-independent);
\item the \textbf{slow} regime (S) yields $\alpha$ via the
  size exponent $W^2\asymp L_0^{2\alpha}t^{2\gamma\alpha}$ at fixed $t$ (equivalently
  via the time exponent $2\gamma\alpha$ with measured $\gamma$).
\end{itemize}
Therefore $\alpha$ and $\beta$ are \emph{unambiguously} extractable, but data from
\emph{both} regimes are required; the control protocol of
Theorem~\ref{thm:control} (tuning $\gamma$ across $1/z$) is exactly the tool that
makes both regimes accessible in one biological system.

\subsection*{Protocol (biological realization)}
\begin{enumerate}
\item Prepare $m\ge2$ replicate colonies/tissues with controlled distinct initial
  front sizes $L_0^{(1)}<\cdots<L_0^{(m)}$, all other conditions identical so that
  $\alpha,\beta,z,f$ are shared.
\item For each replicate and a set of large times $t_1<\cdots<t_n$ (in the
  asymptotic window $t\gg t_0$), image the front, extract $h(x,t)$, and compute the
  spatial variance $W^2(t)=\overline{h^2}-\overline h^2$ (replicate-averaged).
\item Track the front size $L(t)$ and fit $\log L\sim\gamma\log t$ to measure
  $\gamma$.
\item Run the system, via an external control of the expansion rate, in a
  \emph{fast} setting $\gamma_{\mathrm F}>1/z$ and in a \emph{slow} setting
  $\gamma_{\mathrm S}<1/z$. (If only one setting is accessible, only one exponent
  is obtainable; Theorem~\ref{thm:control} gives the minimal change
  $|\delta|>|1/z-\gamma|$ needed to reach the other regime.)
\end{enumerate}

\subsection*{Extraction algorithm}

\begin{algorithm}[H]
\caption{Unambiguous extraction of $\alpha,\beta,z$ from both regimes}
\begin{algorithmic}[1]
\Require Variances $W^2_{\mathrm F}(t_j)$ in a fast setting $\gamma_{\mathrm F}>1/z$;
  variances $W^2_{\mathrm S}(L_0^{(i)},t_j)$ in a slow setting
  $\gamma_{\mathrm S}<1/z$ for $m\ge2$ initial sizes; measured
  $\gamma_{\mathrm F},\gamma_{\mathrm S}$.
\Ensure $\hat\beta,\hat\alpha,\hat z$ with self-consistency checks.
\State \textbf{($\beta$ from fast regime.)} Regress
  $\log W^2_{\mathrm F}(t_j)$ on $\log t_j$; the slope is $2\hat\beta$. Set
  $\hat\beta\gets\tfrac12\cdot\mathrm{slope}$. (Check: this slope is independent of
  $L_0$, per Lemma~\ref{lem:sizeindep}.)
\State \textbf{($\alpha$ from slow regime, size route.)} For each fixed $t_j$,
  regress $\log W^2_{\mathrm S}(L_0^{(i)},t_j)$ on $\log L_0^{(i)}$ across
  $i=1,\dots,m$; the slope is $2\hat\alpha_j$. Check
  $\hat\alpha_1\approx\cdots\approx\hat\alpha_n$ (time-independence of the size
  exponent) and set $\hat\alpha\gets\tfrac12\,\mathrm{mean}_j(2\hat\alpha_j)$.
\State \textbf{($\alpha$ from slow regime, time route — consistency.)} Regress
  $\log W^2_{\mathrm S}$ on $\log t_j$ at fixed $L_0$; the slope is
  $2\gamma_{\mathrm S}\hat\alpha'$, so $\hat\alpha'\gets\mathrm{slope}/(2\gamma_{\mathrm S})$.
  Check $\hat\alpha'\approx\hat\alpha$.
\State \textbf{(Dynamic exponent.)} $\hat z\gets\hat\alpha/\hat\beta$.
\State \textbf{(Regime checks.)} Verify $\gamma_{\mathrm F}>1/\hat z>\gamma_{\mathrm S}$.
\State \Return $\hat\alpha,\hat\beta,\hat z$.
\end{algorithmic}
\end{algorithm}

\begin{proposition}[Correctness, identifiability, termination, complexity]\label{prop:alg}
Under \eqref{eq:FV}--\eqref{eq:adiabatic}:
\begin{enumerate}
\item[\textnormal{(i)}] \emph{(Consistency.)} In the fast setting
  ($\gamma_{\mathrm F}>1/z$), $\log W^2_{\mathrm F}(t)=2\beta\log t+\text{const}+o(\log t)$
  by \eqref{eq:fast}, so the Step-1 slope $\to2\beta$ and $\hat\beta\to\beta$. In
  the slow setting ($\gamma_{\mathrm S}<1/z$),
  $\log W^2_{\mathrm S}(L_0,t)=2\alpha\log L_0+2\gamma_{\mathrm S}\alpha\log t+\text{const}+o(\log t)$
  by \eqref{eq:slow}, so the Step-2 size slope $\to2\alpha$ and the Step-3 time
  slope $\to2\gamma_{\mathrm S}\alpha$; hence $\hat\alpha\to\alpha$,
  $\hat\alpha'\to\alpha$, and $\hat z\to z$.
\item[\textnormal{(ii)}] \emph{(Unambiguous identifiability.)} In the slow regime,
  $(\log L_0,\log t)\mapsto\log W^2_{\mathrm S}$ is affine with gradient
  $(2\alpha,2\gamma_{\mathrm S}\alpha)$. The two regressors $\log L_0$ (varied over
  $m\ge2$ distinct sizes) and $\log t$ (varied over $n\ge2$ times) are
  independently controlled, so the $(\le mn)\times3$ design matrix has full column
  rank and $\alpha$ is uniquely determined; the fast regression gives $\beta$
  separately. Thus $(\alpha,\beta)$ — and $z=\alpha/\beta$ — are uniquely fixed.
\item[\textnormal{(iii)}] \emph{(Why both regimes are needed.)} By
  Lemma~\ref{lem:sizeindep} and \eqref{eq:kappa}, the fast regime is
  $L_0$-independent, so it cannot supply $\alpha$; the slow regime saturates and
  cannot supply $\beta$ (the growth exponent is invisible once the interface is
  saturated). Hence no single regime determines both exponents, and the two-regime
  protocol is necessary; the size route (Step 2) and time route (Step 3) in the
  slow regime cross-check $\alpha$, and Step 6 verifies the regime assignment.
\item[\textnormal{(iv)}] \emph{(Termination/complexity.)} The algorithm performs
  $O(1)$ regressions over $\le mn$ points; it runs in $O(mn)$ time and $O(mn)$
  space and halts (loops over finite index sets only).
\end{enumerate}
\end{proposition}
\begin{proof}
(i) Taking logs of \eqref{eq:fast} and \eqref{eq:slow} gives the affine laws with
$o(\log t)$ residual from the $\asymp$ (the ratio of $W^2$ to its leading power is
bounded between two positive constants and converges, contributing $o(\log t)$ to
the log; ordinary least squares of a function that is affine up to $o(\log t)$
returns slopes converging to the true exponents). (ii) Stack the slow data as
$y=\Phi w$, rows $(1,\log L_0^{(i)},\log t_j)$, $w=(\text{const},2\alpha,
2\gamma_{\mathrm S}\alpha)^\top$; with $m\ge2$ distinct $L_0$ and $n\ge2$ distinct
$t$ the three columns (constant; nonconstant function of $i$; nonconstant function
of $j$) are linearly independent, so $\Phi$ has rank $3$ and $w$ — hence $\alpha$ —
is unique; $\beta$ comes from the rank-$2$ fast fit. (iii) Is
Lemma~\ref{lem:sizeindep} plus \eqref{eq:kappa}: fast amplitude has $L_0$-exponent
$0$, slow width has no $\beta$-dependence (the saturated value $f_\infty L(t)^{2\alpha}$
involves only $\alpha,\gamma$). (iv) Immediate from the operation count over finite
index sets.
\end{proof}

\begin{remark}[Closure of part (b)]
Proposition~\ref{prop:alg} yields $\hat z=\hat\alpha/\hat\beta$, and with the
measured $\gamma$ the order parameter $\widehat\Delta=1/\hat z-\gamma$ of
Proposition~\ref{prop:sign} classifies the phase; by Theorem~\ref{thm:control} the
minimal control strength to switch phase is
$|\delta_{\min}|=|1/\hat z-\gamma|=|\widehat\Delta|$.
\end{remark}

\section{Summary}

\begin{itemize}
\item \textbf{Trichotomy.} Under the adiabatic Family--Vicsek closure, the
  long-time width is the single power law $W^2(t)\asymp\kappa(L_0)\,t^{\theta(\gamma)}$,
  $\theta(\gamma)=\min\{2\beta,2\alpha\gamma\}$ (Theorem~\ref{thm:main}). The two
  announced regimes (F) $\gamma>1/z$ and (S) $\gamma<1/z$ are exactly the two
  leading behaviours, separated by the single critical point $\gamma=1/z$ (regime
  (C)); there are no others. This settles ``prove or disprove'': the two strict
  regimes are the only leading power laws, and the boundary case (omitted in the
  problem) is identified.
\item \textbf{Correction.} The problem's first fast-form $L_0^{2\alpha}t^{2\beta}$
  is inconsistent with FV scaling; the correct fast width is the $L_0$-independent
  $t^{2\beta}$ (Lemma~\ref{lem:sizeindep}, Remark~\ref{rem:correction}). The $L_0$
  dependence resides in the slow regime, $W^2\asymp L_0^{2\alpha}t^{2\gamma\alpha}$.
\item \textbf{Order parameter \& control.} $\Delta=1/z-\gamma$ equals
  $\lim_t\log(\xi/L)/\log t$; its sign is the phase and $|\Delta|$ the per-decade
  rate (Proposition~\ref{prop:sign}). Minimal control to switch phase is
  $|\delta_{\min}|=|1/z-\gamma_0|$ (Theorem~\ref{thm:control}).
\item \textbf{Extraction.} $\beta$ from the fast regime (time slope), $\alpha$ from
  the slow regime (size slope at fixed $t$, cross-checked by the time slope divided
  by $\gamma$); both regimes are necessary and the parameters are identifiable
  (Algorithm 1, Proposition~\ref{prop:alg}).
\end{itemize}

\end{document}
